% ============================================================
\section{Conclusion}
\label{sec:conclusion}

We presented a Gaussian potential-field heuristic for the placement sub-problem of
fault-tolerant compilation. The method superposes attractive magic-state and
partner fields with repulsive fields and places each qubit at the safe maximum of
the result, folding $T$-gate locality, interaction clustering, and corridor
preservation into a single fast scoring pass while safe passage keeps every
mapping routable. Implemented in an open-source OpenQASM-to-lattice-surgery
compiler and evaluated against a randomized-restart baseline and the WISQ
compiler, it \todo{state the headline result once the tables are filled}.

The method has four main limitations. First, placement is static and single-pass:
once a qubit is committed, the mapper does not revise earlier choices in response
to congestion revealed during routing. Second, the field weights are hand-tuned
rather than learned or jointly optimized, which may limit generalization across
workload classes. Third, the heuristic optimizes initial-placement quality and
downstream routing steps, but does not explicitly co-optimize architecture-level
bottlenecks such as dynamic factory throughput, logical-error accumulation, or
event-driven resource contention~\cite{kan2025sparo,bharadwaj2026qomet}. Fourth,
our experimental comparison is deliberately narrow---a randomized baseline and
WISQ---rather than against topology-aware or architecture-aware
compilers~\cite{zhou2026topols,kan2025sparo}. These limitations are consistent
with the broader finding that no single lattice-surgery compilation family is
universally optimal, the preferred strategy depending on circuit density,
logical-qubit count and $T$ fraction~\cite{leblond2025comparison}.

These limitations point to natural extensions: coupling placement to router
feedback so the field adapts online to observed congestion; replacing hand-tuned
coefficients with workload-conditioned or learned parameters; and evaluating the
field placement as an initializer for stronger downstream optimizers---including
topology-aware compilers and architecture-aware
schedulers~\cite{zhou2026topols,bharadwaj2026qomet}---or against emerging
alternatives to static place-and-route such as movable logical
qubits~\cite{herzog2025movable}. More broadly, our results suggest practical value
in lightweight, interpretable front-end mappers: a deterministic field-based
placement occupies a useful middle ground, capturing the key geometric constraints
of lattice-surgery execution at very low computational cost while serving as a
strong standalone baseline and a possible initialization layer for more expensive
optimization flows. Seen this way, the magic-state--aware attractive fields are a
local proxy for the broader class of workload-dependent resource bottlenecks
identified in recent fault-tolerant resource
analyses~\cite{ogorman2017factories,leblond2025comparison}.
